% Part: normal-modal-logic
% Chapter: filtrations
% Section: finite

\documentclass[../../../include/open-logic-section]{subfiles}

\begin{document}

\olfileid{nml}{fil}{fin}

\olsection{Filtrasi Iku Winates}

Kita wis netepake filtrasi kanggo sembarang himpunan~$\Gamma$ kang
katutup tumrap sub-!!{formula}s. Definisi iku dhewe ora njamin yen
filtrasi winates. Malah, yen $\Gamma$ tanpa wates (upamane himpunan
kabeh !!{formula}s), filtrasi bisa tanpa wates. Nanging yen
$\Gamma$ winates (upamane himpunan sub-!!{formula}s saka
!!{formula}~$!A$ tartamtu), saben filtrasi liwat~$\Gamma$ uga
winates.

\begin{prop}\ollabel{prop:filt-are-finite}
  Yen $\Gamma$ winates, saben filtrasi $\mModel{M^*}$ saka model
  $\mModel{M}$ liwat $\Gamma$ uga winates.
\end{prop}

\begin{proof}
  Ukuran $W^*$ yaiku cacah kelas~$[w]$ kang beda miturut relasi
  ekuivalensi~$\equiv$. Sembarang rong jagad $u$, $v$ ing kelas kang
  padha---yaiku sembarang $u$ lan $v$ kang~$u \equiv v$---sarujuk
  tumrap kabeh !!{formula}s~$!A$ ing~$\Gamma$: kanggo $!A \in \Gamma$,
  $!A$ bisa bener ing $u$ lan $v$ loro-lorone, utawa salah ing
  loro-lorone. Dadi saben kelas~$[w]$ cocog karo sawijining himpunan
  bagean saka~$\Gamma$, yaiku himpunan kabeh $!A \in \Gamma$ saengga
  $!A$ bener ing jagad-jagad saka~$[w]$. Ora ana rong kelas beda $[u]$ lan
  $[v]$ kang cocog karo himpunan bagean saka~$\Gamma$ kang padha.
  Amarga yen himpunan !!{formula}s kang bener ing $u$ lan himpunan
  !!{formula}s kang bener ing $v$ padha, mula $u$ lan $v$ sarujuk
  tumrap kabeh formula ing~$\Gamma$, yaiku $u \equiv v$. Nanging
  banjur $[u] = [v]$. Dadi ana fungsi !!a{injective} saka $W^*$
  menyang $\Pow{\Gamma}$, mula
  $\card{W^*} \le \card{\Pow{\Gamma}}$. Mula yen $\Gamma$ ngemot
  $n$ ukara, kardinalitas $W^*$ ora luwih saka~$2^n$.
\end{proof}

\end{document}

